\section{Conclusão}

O arco temporal de maturação das plataformas \textit{open-source} na odontologia digital superou a estagnação metodológica das provas de conceito rudimentares e alcançou a crista de uma mudança paradigmática sem precedentes. O ecossistema multifacetado -- amparado por pilares geométricos, de visualização e análise de elementos finitos como 3D Slicer, MeshLab, Calculix, e FEniCS -- cimenta a autonomia do pesquisador ao estilhaçar as restrições impostas pelos oligopólios de *softwares* proprietários \cite{robles2023opentwins}.

Nessa marcha rumo à medicina hiperpersonalizada e preditiva \cite{digital2026technology}, os modelos paramétricos não subsistem em isolamento. A arquitetura propulsora de plataformas compostas como o OpenCode Ecosystem e as ontologias do OpenTwins inserem a automação profunda de IA diretamente no fluxo da modelagem mecânica. Com seus 128 agentes, escudos governamentais atestados pelo TrustEngine, suítes de validação rigorosas e integrações \textit{Model Context Protocol}, o OpenCode se revela o maquinário ideal para forjar não apenas um modelo computacional da biologia oral, mas um verdadeiro construto epistemológico auditável.

\destaque{A consolidação de gêmeos digitais transcende a adoção de tecnologias reativas e institui uma prática de antecipação preventiva. Em tal cenário, erros iatrogênicos são refutados nas malhas computacionais e jamais nos tecidos biológicos da humanidade.}

As fundações tecnológicas encontram-se indubitavelmente assentadas. No entanto, o ápice desta revolução dependerá da superação inclemente das assimetrias sociodigitais, da validação laboratorial mediante ensaios clínicos rigorosos multicêntricos, e de uma governança ética inflexível e contínua \cite{realising2026digital, frota2025sus}. Esta agenda colaborativa em redes \textit{open-source} oferece, talvez pela primeira vez na história da especialidade, os alicerces necessários para edificar uma odontologia verdadeiramente equitativa, metodologicamente inexpugnável e transparente.
